\chapter{Mạng: Service, DNS, Ingress}
\label{ch:networking}

\section{Vấn đề mà pod tạo ra}

Pod có IP định tuyến được trong cụm --- và mất nó mỗi lần bị thay thế.
Chương~15 đã biến pod thành thứ dùng xong bỏ; phải có gì đó ổn định đứng
phía trước. Thứ đó là \emph{Service}: một IP ảo và tên DNS cố định, phía
sau là danh sách sống các endpoint pod được chọn theo label:

\begin{lstlisting}[language=yaml]
apiVersion: v1
kind: Service
metadata: { name: course-service }
spec:
  selector: { app: course-service }     (*@\co{1}@*)
  ports:
    - { port: 8082, targetPort: 8082 }  (*@\co{2}@*)
\end{lstlisting}

\co{1}~Endpoint controller liên tục liệt kê các pod \emph{Ready} khớp
selector này --- readiness (Chương~17) là bài kiểm tra tư cách thành
viên. \co{2}~Traffic tới port của Service được chuyển tiếp tới targetPort
của một endpoint được chọn.

DNS của cụm cho mỗi Service một cái tên: \code{course-service} trong
namespace, \code{course-service.ts.svc.cluster.local} từ bất cứ đâu.
Đây là lý do các override biến môi trường trên Kubernetes trông giống
hệt bên Compose --- cả hai nền tảng đều phân giải tên service trên một
mạng riêng; dự án cố tình đặt tên Service k8s theo tên service Compose
để các override khớp nhau.

Kiểu Service mặc định, \code{ClusterIP}, chỉ với tới được \emph{từ bên
trong} cụm --- một thuộc tính bảo mật mà Phần~II đã dựa vào cho Mongo
không mật khẩu. \code{NodePort} và \code{LoadBalancer} tồn tại để expose
Service trực tiếp; dự án này thay vào đó dồn mọi thứ qua một Ingress.

\section{Ingress và Traefik}

Một Service expose một thứ; \emph{Ingress} là cửa trước ở tầng L7: định
tuyến HTTP theo host và path, sau này là kết thúc TLS, được khai báo
dưới dạng dữ liệu và thực thi bởi một \emph{Ingress controller} ---
trong k3s là Traefik đi kèm, theo dõi các object Ingress. Từ
\code{api-gateway.yaml}:

\begin{lstlisting}[language=yaml]
kind: Ingress
spec:
  rules:
    - host: ts.local                    (*@\co{1}@*)
      http:
        paths:
          - path: /api                  (*@\co{2}@*)
            pathType: Prefix
            backend:
              service:
                name: api-gateway
                port: { number: 8080 }
          # /oauth2 and /login/oauth2 likewise -> api-gateway
\end{lstlisting}

\begin{description}
  \item[\co{1}] Quy tắc theo host: áp dụng cho request có header
  \code{Host} là \code{ts.local} --- một cái tên chỉ tồn tại trong file
  hosts của Windows, trỏ tới IP tailnet của server. Nhiều host có thể
  dùng chung cặp cổng 80/443 duy nhất của cụm.
  \item[\co{2}] Tỏa theo path. Chỉ ba prefix được vào, tất cả tới
  gateway; \code{/} để trống cho Ingress riêng của frontend
  (Chương~26), mà Traefik gộp vào bảng định tuyến của cùng host.
\end{description}

\section{Đường đi trọn vẹn của một request}

\begin{figure}[h]
\centering
\begin{tikzpicture}[node distance=3.5mm and 5.5mm]
  \node[boxdark] (b) {Browser};
  \node[box, right=of b] (t) {Traefik};
  \node[box, right=of t] (gws) {Service\\\scriptsize api-gateway};
  \node[box, right=of gws] (gw) {pod gateway};
  \node[box, below=of gw] (cs) {Service\\\scriptsize course-service};
  \node[box, left=of cs] (csp) {pod course};
  \draw[flow] (b) -- node[lbl,above]{\scriptsize Host: ts.local} (t);
  \draw[flow] (t) -- node[lbl,above]{\scriptsize /api} (gws);
  \draw[flow] (gws) -- (gw);
  \draw[flow] (gw) -- node[lbl,right]{\scriptsize lb:// qua Eureka}(cs);
  \draw[flow] (cs) -- (csp);
\end{tikzpicture}
\caption{Hai bộ định tuyến, hai hệ thống khám phá, một request.}
\end{figure}

Đáng dừng lại một chút: tra cứu Eureka của gateway trả về \emph{IP pod}
(prefer-ip-address của Chương~3), nên chặng thứ tư về kỹ thuật đi vòng
qua Service course-service. Hai hệ thống khám phá chạy song song ---
Eureka cho gateway-tới-service, DNS của Service cho mọi thứ còn lại
(config-server, cơ sở dữ liệu, Kafka). Bước đơn giản hóa ở Chương~27
gộp chúng thành một.

\begin{pitfall}[Service không chọn được gì]
Triệu chứng: \code{curl} tới tên Service bị timeout; pod đang Running và
khỏe mạnh. Nguyên nhân \#1: lệch label --- selector của Service và label
trong pod template khác nhau bởi một lỗi gõ; danh sách endpoint rỗng và
không ai kêu ca. Kiểm tra: \code{kubectl -n ts get endpoints
course-service} --- cột ENDPOINTS trống chính là lỗi này. Nguyên nhân
\#2: endpoint có nhưng pod chưa Ready --- cùng lệnh đó hiện chúng dưới
\code{NotReadyAddresses}. Lớp trừu tượng Service thất bại
\emph{trong im lặng}; object endpoints là nơi nó lộ ra.
\end{pitfall}

\begin{decision}[một host Ingress + file hosts, không dùng NodePort]
NodePort (mỗi service một cổng lạ lẫm) sẽ chạy được qua Tailscale mà
không cần trò DNS nào --- và tạo ra URL kiểu
\code{100.85.66.43:30082} vừa lộ topology, vừa né gateway, vừa không
làm TLS tử tế được. Một host Ingress nghĩa là một địa chỉ cần bảo vệ và
trải nghiệm trình duyệt giống hệt production; dòng \code{ts.local}
trong file hosts là chi phí duy nhất phía client. Đường nâng cấp ---
một domain thật cộng cert-manager, hoặc Cloudflare Tunnel --- chỉ đổi
host của Ingress và không gì khác.
\end{decision}

\section{Ngoài phạm vi dự án}

Trên EKS, cùng object Ingress đó, thêm annotation cho controller của
AWS, sẽ cấp phát một ALB thật; trên OpenShift nó trở thành
\code{Route} --- cùng ý tưởng, khác danh từ. Mảnh ghép cụm này bỏ qua:
\textbf{NetworkPolicy} --- quy tắc tường lửa cho traffic pod-tới-pod
(``chỉ gateway được chạm tới course-service:8082''), tức là sự thực thi
trong cụm mà mẫu biên-tin-cậy của Chương~5 gọi là điểm mềm còn lại của
nó. Tin cậy mạng phẳng là bình thường với homelab và là cẩu thả với một
ngân hàng.

\begin{quiz}
\begin{enumerate}
  \item IP pod đổi mỗi lần thay thế, vậy mà \code{course-service:8082}
  luôn chạy. Kể tên từng cơ chế nằm giữa cái tên và socket của pod.
  \item Readiness tương tác với endpoints của Service ra sao, và điều
  đó nghĩa là gì trong một rolling update?
  \item Chẩn đoán: Service tồn tại, pod Ready, \code{get endpoints}
  hiện danh sách rỗng. Sai ở đâu, và vì sao không có gì báo lỗi?
  \item Vì sao \code{ts.local} cần một dòng trong file hosts, và thứ gì
  thay thế dòng đó trong hai đường nâng cấp lên production?
\end{enumerate}
\end{quiz}
